\section{Práctica 2:}

\subsection{Función lógica por SOP y POS}

En la tabla~\ref{tab:truth_table_min_max} se muestra la tabla de verdad tanto para la función booleana de mintérminos (ecuación~\ref{eq:booleana_min}) como para maxitérminos ( ecuación~\ref{eq:booleana_max}). Luego se procede a reducir dichas funciones con el mapa de Karnaugh para hallar el circuito combinacional reducido para cada salida.

% Please add the following required packages to your document preamble:
% \usepackage{multirow}
% \usepackage[table,xcdraw]{xcolor}
% Beamer presentation requires \usepackage{colortbl} instead of \usepackage[table,xcdraw]{xcolor}
\begin{table}
\centering
\caption{Tabla de verdad para funciones booleanas (min. y max.).}
\label{tab:truth_table_min_max}
\begin{tabular}{|
>{\columncolor[HTML]{FFFFFF}}c |
>{\columncolor[HTML]{FFFFFF}}c 
>{\columncolor[HTML]{FFFFFF}}c 
>{\columncolor[HTML]{FFFFFF}}c 
>{\columncolor[HTML]{FFFFFF}}c |
>{\columncolor[HTML]{FFFFFF}}c 
>{\columncolor[HTML]{FFFFFF}}c |}
\hline
\cellcolor[HTML]{FFFFFF}                     & \multicolumn{4}{c|}{\cellcolor[HTML]{FFFFFF}\textbf{Entradas}}                                                                                                                           & \multicolumn{2}{c|}{\cellcolor[HTML]{FFFFFF}\textbf{Salida}}                                                                                     \\ \cline{2-7} 
\multirow{-2}{*}{\cellcolor[HTML]{FFFFFF}\#} & \multicolumn{1}{c|}{\cellcolor[HTML]{FFFFFF}\textbf{A}} & \multicolumn{1}{c|}{\cellcolor[HTML]{FFFFFF}\textbf{B}} & \multicolumn{1}{c|}{\cellcolor[HTML]{FFFFFF}\textbf{C}} & \textbf{D} & \multicolumn{1}{c|}{\cellcolor[HTML]{FFFFC7}{\color[HTML]{013300} \textbf{Min.}}} & \cellcolor[HTML]{ECF4FF}{\color[HTML]{00009B} \textbf{Max.}} \\ \hline
{\color[HTML]{9B9B9B} 0}                     & \multicolumn{1}{c|}{\cellcolor[HTML]{FFFFFF}0}          & \multicolumn{1}{c|}{\cellcolor[HTML]{FFFFFF}0}          & \multicolumn{1}{c|}{\cellcolor[HTML]{FFFFFF}0}          & 0          & \multicolumn{1}{c|}{\cellcolor[HTML]{FFFFFF}0}                                    & 1                                                            \\ \hline
{\color[HTML]{9B9B9B} 1}                     & \multicolumn{1}{c|}{\cellcolor[HTML]{FFFFFF}0}          & \multicolumn{1}{c|}{\cellcolor[HTML]{FFFFFF}0}          & \multicolumn{1}{c|}{\cellcolor[HTML]{FFFFFF}0}          & 1          & \multicolumn{1}{c|}{\cellcolor[HTML]{FFFFC7}{\color[HTML]{013300} \textbf{1}}}    & {\color[HTML]{333333} 1}                                     \\ \hline
{\color[HTML]{9B9B9B} 2}                     & \multicolumn{1}{c|}{\cellcolor[HTML]{FFFFFF}0}          & \multicolumn{1}{c|}{\cellcolor[HTML]{FFFFFF}0}          & \multicolumn{1}{c|}{\cellcolor[HTML]{FFFFFF}1}          & 0          & \multicolumn{1}{c|}{\cellcolor[HTML]{FFFFFF}0}                                    & 1                                                            \\ \hline
{\color[HTML]{9B9B9B} 3}                     & \multicolumn{1}{c|}{\cellcolor[HTML]{FFFFFF}0}          & \multicolumn{1}{c|}{\cellcolor[HTML]{FFFFFF}0}          & \multicolumn{1}{c|}{\cellcolor[HTML]{FFFFFF}1}          & 1          & \multicolumn{1}{c|}{\cellcolor[HTML]{FFFFC7}{\color[HTML]{013300} \textbf{1}}}    & \cellcolor[HTML]{ECF4FF}{\color[HTML]{00009B} \textbf{0}}    \\ \hline
{\color[HTML]{9B9B9B} 4}                     & \multicolumn{1}{c|}{\cellcolor[HTML]{FFFFFF}0}          & \multicolumn{1}{c|}{\cellcolor[HTML]{FFFFFF}1}          & \multicolumn{1}{c|}{\cellcolor[HTML]{FFFFFF}0}          & 0          & \multicolumn{1}{c|}{\cellcolor[HTML]{FFFFFF}0}                                    & 1                                                            \\ \hline
{\color[HTML]{9B9B9B} 5}                     & \multicolumn{1}{c|}{\cellcolor[HTML]{FFFFFF}0}          & \multicolumn{1}{c|}{\cellcolor[HTML]{FFFFFF}1}          & \multicolumn{1}{c|}{\cellcolor[HTML]{FFFFFF}0}          & 1          & \multicolumn{1}{c|}{\cellcolor[HTML]{FFFFC7}{\color[HTML]{013300} \textbf{1}}}    & \cellcolor[HTML]{ECF4FF}{\color[HTML]{00009B} \textbf{0}}    \\ \hline
{\color[HTML]{9B9B9B} 6}                     & \multicolumn{1}{c|}{\cellcolor[HTML]{FFFFFF}0}          & \multicolumn{1}{c|}{\cellcolor[HTML]{FFFFFF}1}          & \multicolumn{1}{c|}{\cellcolor[HTML]{FFFFFF}1}          & 0          & \multicolumn{1}{c|}{\cellcolor[HTML]{FFFFFF}0}                                    & 1                                                            \\ \hline
{\color[HTML]{9B9B9B} 7}                     & \multicolumn{1}{c|}{\cellcolor[HTML]{FFFFFF}0}          & \multicolumn{1}{c|}{\cellcolor[HTML]{FFFFFF}1}          & \multicolumn{1}{c|}{\cellcolor[HTML]{FFFFFF}1}          & 1          & \multicolumn{1}{c|}{\cellcolor[HTML]{FFFFC7}{\color[HTML]{013300} \textbf{1}}}    & \cellcolor[HTML]{ECF4FF}{\color[HTML]{00009B} \textbf{0}}    \\ \hline
{\color[HTML]{9B9B9B} 8}                     & \multicolumn{1}{c|}{\cellcolor[HTML]{FFFFFF}1}          & \multicolumn{1}{c|}{\cellcolor[HTML]{FFFFFF}0}          & \multicolumn{1}{c|}{\cellcolor[HTML]{FFFFFF}0}          & 0          & \multicolumn{1}{c|}{\cellcolor[HTML]{FFFFFF}0}                                    & 1                                                            \\ \hline
{\color[HTML]{9B9B9B} 9}                     & \multicolumn{1}{c|}{\cellcolor[HTML]{FFFFFF}1}          & \multicolumn{1}{c|}{\cellcolor[HTML]{FFFFFF}0}          & \multicolumn{1}{c|}{\cellcolor[HTML]{FFFFFF}0}          & 1          & \multicolumn{1}{c|}{\cellcolor[HTML]{FFFFC7}{\color[HTML]{013300} \textbf{1}}}    & {\color[HTML]{333333} 1}                                     \\ \hline
{\color[HTML]{9B9B9B} 10}                    & \multicolumn{1}{c|}{\cellcolor[HTML]{FFFFFF}1}          & \multicolumn{1}{c|}{\cellcolor[HTML]{FFFFFF}0}          & \multicolumn{1}{c|}{\cellcolor[HTML]{FFFFFF}1}          & 0          & \multicolumn{1}{c|}{\cellcolor[HTML]{FFFFFF}0}                                    & 1                                                            \\ \hline
{\color[HTML]{9B9B9B} 11}                    & \multicolumn{1}{c|}{\cellcolor[HTML]{FFFFFF}1}          & \multicolumn{1}{c|}{\cellcolor[HTML]{FFFFFF}0}          & \multicolumn{1}{c|}{\cellcolor[HTML]{FFFFFF}1}          & 1          & \multicolumn{1}{c|}{\cellcolor[HTML]{FFFFC7}{\color[HTML]{013300} \textbf{1}}}    & \cellcolor[HTML]{ECF4FF}{\color[HTML]{00009B} \textbf{0}}    \\ \hline
{\color[HTML]{9B9B9B} 12}                    & \multicolumn{1}{c|}{\cellcolor[HTML]{FFFFFF}1}          & \multicolumn{1}{c|}{\cellcolor[HTML]{FFFFFF}1}          & \multicolumn{1}{c|}{\cellcolor[HTML]{FFFFFF}0}          & 0          & \multicolumn{1}{c|}{\cellcolor[HTML]{FFFFFF}0}                                    & 1                                                            \\ \hline
{\color[HTML]{9B9B9B} 13}                    & \multicolumn{1}{c|}{\cellcolor[HTML]{FFFFFF}1}          & \multicolumn{1}{c|}{\cellcolor[HTML]{FFFFFF}1}          & \multicolumn{1}{c|}{\cellcolor[HTML]{FFFFFF}0}          & 1          & \multicolumn{1}{c|}{\cellcolor[HTML]{FFFFC7}{\color[HTML]{013300} \textbf{1}}}    & \cellcolor[HTML]{ECF4FF}{\color[HTML]{00009B} \textbf{0}}    \\ \hline
{\color[HTML]{9B9B9B} 14}                    & \multicolumn{1}{c|}{\cellcolor[HTML]{FFFFFF}1}          & \multicolumn{1}{c|}{\cellcolor[HTML]{FFFFFF}1}          & \multicolumn{1}{c|}{\cellcolor[HTML]{FFFFFF}1}          & 0          & \multicolumn{1}{c|}{\cellcolor[HTML]{FFFFFF}0}                                    & 1                                                            \\ \hline
{\color[HTML]{9B9B9B} 15}                    & \multicolumn{1}{c|}{\cellcolor[HTML]{FFFFFF}1}          & \multicolumn{1}{c|}{\cellcolor[HTML]{FFFFFF}1}          & \multicolumn{1}{c|}{\cellcolor[HTML]{FFFFFF}1}          & 1          & \multicolumn{1}{c|}{\cellcolor[HTML]{FFFFC7}{\color[HTML]{013300} \textbf{1}}}    & \cellcolor[HTML]{ECF4FF}{\color[HTML]{00009B} \textbf{0}}    \\ \hline
\end{tabular}
\end{table}



\vspace{2.0mm}
\subsubsection{Mintérminos}

\begin{equation}
    F \left(A,B,B,C,D \right) = \sum{\left ( 1,3,5,7,9,11,13,15 \right )}
    \label{eq:booleana_min}
\end{equation}

De la tabla de verdad (tabla~\ref{tab:truth_table_min_max}) se tienen en cuenta los $1$ para la salida por mintérminos. La reducción utilizando Karnaugh se muestra en la figura~\ref{fig:Karnaugh_min} donde sólo se tiene un grupo de $1$s haciendo que la salida sea:

\begin{equation}
    Salida = D
    \label{eq:out_min}
\end{equation}

\begin{figure}[H]
    \centering
    \begin{Karnaugh}
        \contingut{0,1,0,1,0,1,0,1,0,1,0,1,0,1,0,1}
        \implicant{1}{11}{green}
    \end{Karnaugh}
    \caption{Mapa de Karnaugh para la reducción por mintérminos}.
    \label{fig:Karnaugh_min}
\end{figure}

La comprobación de la reducción se simuló en logisim (ver figura~\ref{fig:logisim_min}), mientras que la implementación se muestra en la figura~\ref{fig:montaje_min}. Para este circuito, la salida se conecta directamente a la entrada $D$. 
Por lo tanto, los estados de las entradas $A$, $B$ y $C$ no importan ($X$) para determinar el estado de la salida. En otras palabras, cuando $D=1$ la salida es $1$ y cuando $D=0$ la salida es $0$.


\begin{figure}[]
    \centering
    \includegraphics[width=0.95\linewidth]{img/sim/3_min_logisim.png}
    \caption{Comprobación de la reducción por suma de productos con logisim.}
    \label{fig:logisim_min}
\end{figure}


\begin{figure}
    \centering
    \includegraphics[width=1.0\linewidth]{img/lab/min_montaje.jpeg}
    \caption{Montaje del circuito por mintérminos.}
    \label{fig:montaje_min}
\end{figure}

%%%%%%%%%%%%%%%%%%%%%%%%%%%%%%%%%%%%%%%%%%
\vspace{2.0mm}
\subsubsection{Maxitérminos}

\begin{equation}
    F \left(A,B,B,C,D \right) = \prod{\left ( 3,5,7,11,13,15) \right )}
    \label{eq:booleana_max}
\end{equation}

De la tabla de verdad (tabla~\ref{tab:truth_table_min_max}) se tienen en cuenta los $0$ para reducir la salida por maxitérminos. La reducción utilizando Karnaugh se muestra en la figura~\ref{fig:Karnaugh_max} donde se tienen dos grupos de $0$ haciendo que la salida sea:

\begin{equation}
    Salida = \left ( {\color{blue} \overline{B} + \overline{D}} \right ) \left ( {\color{orange} \overline{C} + \overline{D}} \right )
    \label{eq:out_max}
\end{equation}

\begin{figure}[]
    \centering
    \begin{Karnaugh}
        \contingut{1,1,1,0,1,0,1,0,1,1,1,0,1,0,1,0}
        \implicant{3}{11}{orange}
        \implicant{5}{15}{blue}
    \end{Karnaugh}
    \caption{Mapa de Karnaugh para la reducción por maxitérminos}.
    \label{fig:Karnaugh_max}
\end{figure}

La comprobación de la reducción se simuló en logisim (ver figura~\ref{fig:logisim_max}). Adicionalmente, utilizando EdaPlayground se sintetizó el hardware para determinar el funcionamiento esperado del circuito. Finalmente, se comprobó la tabla de verdad con el montaje en el laboratorio como se muestra en la figura~\ref{fig:montaje_max} y que se puede consultar en el \href{https://youtu.be/wEyro38Z5Xk}{link del video} alojado en youtube.

\begin{figure}[]
    \centering
    \includegraphics[width=1.0\linewidth]{img/sim/3_max_logisim.png}
    \caption{Comprobación de la reducción por producto de sumas con logisim.}
    \label{fig:logisim_max}
\end{figure}

\begin{figure}[]
    \centering
    \includegraphics[width=1.0\linewidth]{img/sim/maxi_qucs.png}
    \caption{Comprobación de la tabla de verdad con Qucs.}
    \label{fig:qucs_max}
\end{figure}

\begin{figure}
    \centering
    \includegraphics[width=1.0\linewidth]{img/sim/3_max_ise_sim.png}
    \caption{Sintetización de Hardware en ISE de Xilinx.}
    \label{fig:vhdl_max}
\end{figure}

\begin{figure}[H]
    \centering
    \includegraphics[width=1.0\linewidth]{img/lab/maxi_montaje.png}
    \caption{Montaje del circuito por maxitérminos. Comprobación de la tabla de verdad en el \href{https://youtu.be/wEyro38Z5Xk}{link del video}.}
    \label{fig:montaje_max}
\end{figure}

%%%%%%%%%%%%%%%%%%%%%%%%%%%%%%%%%%%%%%%%%%%%%%%%%%%%%%%%%%%

\subsection{Decodificador BCD a 7 segmentos}

Para el decodificador de BCD a 7 segmentos se tiene su tabla de verdad (ver tabla~\ref{tab:truth_table_BCD_a_7_seg}),la simulación en logisim (ver figura~\ref{fig:BCD_7seg_logisim} con su \href{https://youtu.be/_nBrKwiCBEE}{video de simulación}), simulación en proteus (ver figura~\ref{fig:BCD_7seg_proteus} con su \href{https://youtu.be/dHLRmH1Clhk?si=OIODbV-4adk1f7gQ}{video de simulación}), sintetización de hardware en ISE (ver figura~\ref{fig:BCD_7seg_ise}) y el montaje de la práctica (ver figura~\ref{fig:BCD_7seg_montaje} con su \href{https://youtu.be/YC1jDNmOjXA}{video de funcionamiento}).

% Please add the following required packages to your document preamble:
% \usepackage{multirow}
% \usepackage[table,xcdraw]{xcolor}
% Beamer presentation requires \usepackage{colortbl} instead of \usepackage[table,xcdraw]{xcolor}
\begin{table}[t]
\centering
\caption{Tabla de verdad para decodificador de BCD a 7 segmentos (7448).}
\label{tab:truth_table_BCD_a_7_seg}
\begin{tabular}{|
>{\columncolor[HTML]{FFFFFF}}c |
>{\columncolor[HTML]{FFFFFF}}c 
>{\columncolor[HTML]{FFFFFF}}c 
>{\columncolor[HTML]{FFFFFF}}c 
>{\columncolor[HTML]{FFFFFF}}c |
>{\columncolor[HTML]{ECF4FF}}l 
>{\columncolor[HTML]{ECF4FF}}l 
>{\columncolor[HTML]{ECF4FF}}l 
>{\columncolor[HTML]{ECF4FF}}l 
>{\columncolor[HTML]{ECF4FF}}l 
>{\columncolor[HTML]{ECF4FF}}l 
>{\columncolor[HTML]{ECF4FF}}l |}
\hline
\cellcolor[HTML]{FFFFFF}                     & \multicolumn{4}{c|}{\cellcolor[HTML]{FFFFFF}\textbf{BCD}}                                                                                                                                & \multicolumn{7}{c|}{\cellcolor[HTML]{ECF4FF}{\color[HTML]{000000} \textbf{7 Segmentos}}}  \\ \cline{2-12} 
\multirow{-2}{*}{\cellcolor[HTML]{FFFFFF}\#} & \multicolumn{1}{c|}{\cellcolor[HTML]{FFFFFF}\textbf{D}} & \multicolumn{1}{c|}{\cellcolor[HTML]{FFFFFF}\textbf{C}} & \multicolumn{1}{c|}{\cellcolor[HTML]{FFFFFF}\textbf{B}} & \textbf{A} & \multicolumn{1}{c|}{\cellcolor[HTML]{ECF4FF}{\color[HTML]{000000} \textbf{a}}} & \multicolumn{1}{c|}{\cellcolor[HTML]{ECF4FF}{\color[HTML]{000000} \textbf{b}}} & \multicolumn{1}{l|}{\cellcolor[HTML]{ECF4FF}{\color[HTML]{000000} \textbf{c}}} & \multicolumn{1}{l|}{\cellcolor[HTML]{ECF4FF}{\color[HTML]{000000} \textbf{d}}} & \multicolumn{1}{l|}{\cellcolor[HTML]{ECF4FF}{\color[HTML]{000000} \textbf{e}}} & \multicolumn{1}{l|}{\cellcolor[HTML]{ECF4FF}{\color[HTML]{000000} \textbf{f}}} & {\color[HTML]{000000} \textbf{g}}                \\ \hline
{\color[HTML]{9B9B9B} 0}                     & \multicolumn{1}{c|}{\cellcolor[HTML]{FFFFFF}0}          & \multicolumn{1}{c|}{\cellcolor[HTML]{FFFFFF}0}          & \multicolumn{1}{c|}{\cellcolor[HTML]{FFFFFF}0}          & 0          & \multicolumn{1}{c|}{\cellcolor[HTML]{ECF4FF}{\color[HTML]{000000} 1}}          & \multicolumn{1}{c|}{\cellcolor[HTML]{ECF4FF}{\color[HTML]{000000} 1}}          & \multicolumn{1}{l|}{\cellcolor[HTML]{ECF4FF}{\color[HTML]{000000} 1}}          & \multicolumn{1}{l|}{\cellcolor[HTML]{ECF4FF}{\color[HTML]{000000} 1}}          & \multicolumn{1}{l|}{\cellcolor[HTML]{ECF4FF}{\color[HTML]{000000} 1}}          & \multicolumn{1}{l|}{\cellcolor[HTML]{ECF4FF}{\color[HTML]{000000} 1}}          & {\color[HTML]{000000} 0}                         \\ \hline
{\color[HTML]{9B9B9B} 1}                     & \multicolumn{1}{c|}{\cellcolor[HTML]{FFFFFF}0}          & \multicolumn{1}{c|}{\cellcolor[HTML]{FFFFFF}0}          & \multicolumn{1}{c|}{\cellcolor[HTML]{FFFFFF}0}          & 1          & \multicolumn{1}{c|}{\cellcolor[HTML]{ECF4FF}{\color[HTML]{000000} 0}}          & \multicolumn{1}{c|}{\cellcolor[HTML]{ECF4FF}{\color[HTML]{000000} 1}}          & \multicolumn{1}{l|}{\cellcolor[HTML]{ECF4FF}{\color[HTML]{000000} 1}}          & \multicolumn{1}{l|}{\cellcolor[HTML]{ECF4FF}{\color[HTML]{000000} 0}}          & \multicolumn{1}{l|}{\cellcolor[HTML]{ECF4FF}{\color[HTML]{000000} 0}}          & \multicolumn{1}{l|}{\cellcolor[HTML]{ECF4FF}{\color[HTML]{000000} 0}}          & {\color[HTML]{000000} 0}                         \\ \hline
{\color[HTML]{9B9B9B} 2}                     & \multicolumn{1}{c|}{\cellcolor[HTML]{FFFFFF}0}          & \multicolumn{1}{c|}{\cellcolor[HTML]{FFFFFF}0}          & \multicolumn{1}{c|}{\cellcolor[HTML]{FFFFFF}1}          & 0          & \multicolumn{1}{c|}{\cellcolor[HTML]{ECF4FF}{\color[HTML]{000000} 1}}          & \multicolumn{1}{c|}{\cellcolor[HTML]{ECF4FF}{\color[HTML]{000000} 1}}          & \multicolumn{1}{l|}{\cellcolor[HTML]{ECF4FF}{\color[HTML]{000000} 0}}          & \multicolumn{1}{l|}{\cellcolor[HTML]{ECF4FF}{\color[HTML]{000000} 1}}          & \multicolumn{1}{l|}{\cellcolor[HTML]{ECF4FF}{\color[HTML]{000000} 1}}          & \multicolumn{1}{l|}{\cellcolor[HTML]{ECF4FF}{\color[HTML]{000000} 0}}          & {\color[HTML]{000000} 1}                         \\ \hline
{\color[HTML]{9B9B9B} 3}                     & \multicolumn{1}{c|}{\cellcolor[HTML]{FFFFFF}0}          & \multicolumn{1}{c|}{\cellcolor[HTML]{FFFFFF}0}          & \multicolumn{1}{c|}{\cellcolor[HTML]{FFFFFF}1}          & 1          & \multicolumn{1}{c|}{\cellcolor[HTML]{ECF4FF}{\color[HTML]{000000} 1}}          & \multicolumn{1}{c|}{\cellcolor[HTML]{ECF4FF}{\color[HTML]{000000} 1}}          & \multicolumn{1}{l|}{\cellcolor[HTML]{ECF4FF}{\color[HTML]{000000} 1}}          & \multicolumn{1}{l|}{\cellcolor[HTML]{ECF4FF}{\color[HTML]{000000} 1}}          & \multicolumn{1}{l|}{\cellcolor[HTML]{ECF4FF}{\color[HTML]{000000} 0}}          & \multicolumn{1}{l|}{\cellcolor[HTML]{ECF4FF}{\color[HTML]{000000} 0}}          & {\color[HTML]{000000} 1}                         \\ \hline
{\color[HTML]{9B9B9B} 4}                     & \multicolumn{1}{c|}{\cellcolor[HTML]{FFFFFF}0}          & \multicolumn{1}{c|}{\cellcolor[HTML]{FFFFFF}1}          & \multicolumn{1}{c|}{\cellcolor[HTML]{FFFFFF}0}          & 0          & \multicolumn{1}{l|}{\cellcolor[HTML]{ECF4FF}{\color[HTML]{000000} 0}}          & \multicolumn{1}{l|}{\cellcolor[HTML]{ECF4FF}{\color[HTML]{000000} 1}}          & \multicolumn{1}{l|}{\cellcolor[HTML]{ECF4FF}{\color[HTML]{000000} 1}}          & \multicolumn{1}{l|}{\cellcolor[HTML]{ECF4FF}{\color[HTML]{000000} 0}}          & \multicolumn{1}{l|}{\cellcolor[HTML]{ECF4FF}{\color[HTML]{000000} 0}}          & \multicolumn{1}{l|}{\cellcolor[HTML]{ECF4FF}{\color[HTML]{000000} 1}}          & {\color[HTML]{000000} 1}                         \\ \hline
{\color[HTML]{9B9B9B} 5}                     & \multicolumn{1}{c|}{\cellcolor[HTML]{FFFFFF}0}          & \multicolumn{1}{c|}{\cellcolor[HTML]{FFFFFF}1}          & \multicolumn{1}{c|}{\cellcolor[HTML]{FFFFFF}0}          & 1          & \multicolumn{1}{l|}{\cellcolor[HTML]{ECF4FF}{\color[HTML]{000000} 1}}          & \multicolumn{1}{l|}{\cellcolor[HTML]{ECF4FF}{\color[HTML]{000000} 0}}          & \multicolumn{1}{l|}{\cellcolor[HTML]{ECF4FF}{\color[HTML]{000000} 1}}          & \multicolumn{1}{l|}{\cellcolor[HTML]{ECF4FF}{\color[HTML]{000000} 1}}          & \multicolumn{1}{l|}{\cellcolor[HTML]{ECF4FF}{\color[HTML]{000000} 0}}          & \multicolumn{1}{l|}{\cellcolor[HTML]{ECF4FF}{\color[HTML]{000000} 1}}          & {\color[HTML]{000000} 1}                         \\ \hline
{\color[HTML]{9B9B9B} 6}                     & \multicolumn{1}{c|}{\cellcolor[HTML]{FFFFFF}0}          & \multicolumn{1}{c|}{\cellcolor[HTML]{FFFFFF}1}          & \multicolumn{1}{c|}{\cellcolor[HTML]{FFFFFF}1}          & 0          & \multicolumn{1}{l|}{\cellcolor[HTML]{ECF4FF}{\color[HTML]{000000} 0}}          & \multicolumn{1}{l|}{\cellcolor[HTML]{ECF4FF}{\color[HTML]{000000} 0}}          & \multicolumn{1}{l|}{\cellcolor[HTML]{ECF4FF}{\color[HTML]{000000} 1}}          & \multicolumn{1}{l|}{\cellcolor[HTML]{ECF4FF}{\color[HTML]{000000} 1}}          & \multicolumn{1}{l|}{\cellcolor[HTML]{ECF4FF}{\color[HTML]{000000} 1}}          & \multicolumn{1}{l|}{\cellcolor[HTML]{ECF4FF}{\color[HTML]{000000} 1}}          & {\color[HTML]{000000} 1}                         \\ \hline
{\color[HTML]{9B9B9B} 7}                     & \multicolumn{1}{c|}{\cellcolor[HTML]{FFFFFF}0}          & \multicolumn{1}{c|}{\cellcolor[HTML]{FFFFFF}1}          & \multicolumn{1}{c|}{\cellcolor[HTML]{FFFFFF}1}          & 1          & \multicolumn{1}{l|}{\cellcolor[HTML]{ECF4FF}{\color[HTML]{000000} 1}}          & \multicolumn{1}{l|}{\cellcolor[HTML]{ECF4FF}{\color[HTML]{000000} 1}}          & \multicolumn{1}{l|}{\cellcolor[HTML]{ECF4FF}{\color[HTML]{000000} 1}}          & \multicolumn{1}{l|}{\cellcolor[HTML]{ECF4FF}{\color[HTML]{000000} 0}}          & \multicolumn{1}{l|}{\cellcolor[HTML]{ECF4FF}{\color[HTML]{000000} 0}}          & \multicolumn{1}{l|}{\cellcolor[HTML]{ECF4FF}{\color[HTML]{000000} 0}}          & {\color[HTML]{000000} 0}                         \\ \hline
{\color[HTML]{9B9B9B} 8}                     & \multicolumn{1}{c|}{\cellcolor[HTML]{FFFFFF}1}          & \multicolumn{1}{c|}{\cellcolor[HTML]{FFFFFF}0}          & \multicolumn{1}{c|}{\cellcolor[HTML]{FFFFFF}0}          & 0          & \multicolumn{1}{l|}{\cellcolor[HTML]{ECF4FF}{\color[HTML]{000000} 1}}          & \multicolumn{1}{l|}{\cellcolor[HTML]{ECF4FF}{\color[HTML]{000000} 1}}          & \multicolumn{1}{l|}{\cellcolor[HTML]{ECF4FF}{\color[HTML]{000000} 1}}          & \multicolumn{1}{l|}{\cellcolor[HTML]{ECF4FF}{\color[HTML]{000000} 1}}          & \multicolumn{1}{l|}{\cellcolor[HTML]{ECF4FF}{\color[HTML]{000000} 1}}          & \multicolumn{1}{l|}{\cellcolor[HTML]{ECF4FF}{\color[HTML]{000000} 1}}          & {\color[HTML]{000000} 1}                         \\ \hline
{\color[HTML]{9B9B9B} 9}                     & \multicolumn{1}{c|}{\cellcolor[HTML]{FFFFFF}1}          & \multicolumn{1}{c|}{\cellcolor[HTML]{FFFFFF}0}          & \multicolumn{1}{c|}{\cellcolor[HTML]{FFFFFF}0}          & 1          & \multicolumn{1}{l|}{\cellcolor[HTML]{ECF4FF}{\color[HTML]{000000} 1}}          & \multicolumn{1}{l|}{\cellcolor[HTML]{ECF4FF}{\color[HTML]{000000} 1}}          & \multicolumn{1}{l|}{\cellcolor[HTML]{ECF4FF}{\color[HTML]{000000} 1}}          & \multicolumn{1}{l|}{\cellcolor[HTML]{ECF4FF}{\color[HTML]{000000} 0}}          & \multicolumn{1}{l|}{\cellcolor[HTML]{ECF4FF}{\color[HTML]{000000} 0}}          & \multicolumn{1}{l|}{\cellcolor[HTML]{ECF4FF}{\color[HTML]{000000} 1}}          & {\color[HTML]{000000} 1}                         \\ \hline
{\color[HTML]{9B9B9B} 10}                    & \multicolumn{1}{c|}{\cellcolor[HTML]{FFFFFF}1}          & \multicolumn{1}{c|}{\cellcolor[HTML]{FFFFFF}0}          & \multicolumn{1}{c|}{\cellcolor[HTML]{FFFFFF}1}          & 0          & \multicolumn{1}{l|}{\cellcolor[HTML]{ECF4FF}{\color[HTML]{000000} 0}}          & \multicolumn{1}{l|}{\cellcolor[HTML]{ECF4FF}{\color[HTML]{000000} 0}}          & \multicolumn{1}{l|}{\cellcolor[HTML]{ECF4FF}{\color[HTML]{000000} 0}}          & \multicolumn{1}{l|}{\cellcolor[HTML]{ECF4FF}{\color[HTML]{000000} 1}}          & \multicolumn{1}{l|}{\cellcolor[HTML]{ECF4FF}{\color[HTML]{000000} 1}}          & \multicolumn{1}{l|}{\cellcolor[HTML]{ECF4FF}{\color[HTML]{000000} 0}}          & {\color[HTML]{000000} 1}                         \\ \hline
{\color[HTML]{9B9B9B} 11}                    & \multicolumn{1}{c|}{\cellcolor[HTML]{FFFFFF}1}          & \multicolumn{1}{c|}{\cellcolor[HTML]{FFFFFF}0}          & \multicolumn{1}{c|}{\cellcolor[HTML]{FFFFFF}1}          & 1          & \multicolumn{1}{l|}{\cellcolor[HTML]{ECF4FF}{\color[HTML]{000000} 0}}          & \multicolumn{1}{l|}{\cellcolor[HTML]{ECF4FF}{\color[HTML]{000000} 0}}          & \multicolumn{1}{l|}{\cellcolor[HTML]{ECF4FF}{\color[HTML]{000000} 1}}          & \multicolumn{1}{l|}{\cellcolor[HTML]{ECF4FF}{\color[HTML]{000000} 1}}          & \multicolumn{1}{l|}{\cellcolor[HTML]{ECF4FF}{\color[HTML]{000000} 0}}          & \multicolumn{1}{l|}{\cellcolor[HTML]{ECF4FF}{\color[HTML]{000000} 0}}          & {\color[HTML]{000000} 1}                         \\ \hline
{\color[HTML]{9B9B9B} 12}                    & \multicolumn{1}{c|}{\cellcolor[HTML]{FFFFFF}1}          & \multicolumn{1}{c|}{\cellcolor[HTML]{FFFFFF}1}          & \multicolumn{1}{c|}{\cellcolor[HTML]{FFFFFF}0}          & 0          & \multicolumn{1}{l|}{\cellcolor[HTML]{ECF4FF}{\color[HTML]{000000} 0}}          & \multicolumn{1}{l|}{\cellcolor[HTML]{ECF4FF}{\color[HTML]{000000} 1}}          & \multicolumn{1}{l|}{\cellcolor[HTML]{ECF4FF}{\color[HTML]{000000} 0}}          & \multicolumn{1}{l|}{\cellcolor[HTML]{ECF4FF}{\color[HTML]{000000} 0}}          & \multicolumn{1}{l|}{\cellcolor[HTML]{ECF4FF}{\color[HTML]{000000} 0}}          & \multicolumn{1}{l|}{\cellcolor[HTML]{ECF4FF}{\color[HTML]{000000} 1}}          & \cellcolor[HTML]{ECF4FF}{\color[HTML]{000000} 1} \\ \hline
{\color[HTML]{9B9B9B} 13}                    & \multicolumn{1}{c|}{\cellcolor[HTML]{FFFFFF}1}          & \multicolumn{1}{c|}{\cellcolor[HTML]{FFFFFF}1}          & \multicolumn{1}{c|}{\cellcolor[HTML]{FFFFFF}0}          & 1          & \multicolumn{1}{l|}{\cellcolor[HTML]{ECF4FF}{\color[HTML]{000000} 1}}          & \multicolumn{1}{l|}{\cellcolor[HTML]{ECF4FF}{\color[HTML]{000000} 0}}          & \multicolumn{1}{l|}{\cellcolor[HTML]{ECF4FF}{\color[HTML]{000000} 0}}          & \multicolumn{1}{l|}{\cellcolor[HTML]{ECF4FF}{\color[HTML]{000000} 1}}          & \multicolumn{1}{l|}{\cellcolor[HTML]{ECF4FF}{\color[HTML]{000000} 0}}          & \multicolumn{1}{l|}{\cellcolor[HTML]{ECF4FF}{\color[HTML]{000000} 1}}          & \cellcolor[HTML]{ECF4FF}{\color[HTML]{000000} 1} \\ \hline
{\color[HTML]{9B9B9B} 14}                    & \multicolumn{1}{c|}{\cellcolor[HTML]{FFFFFF}1}          & \multicolumn{1}{c|}{\cellcolor[HTML]{FFFFFF}1}          & \multicolumn{1}{c|}{\cellcolor[HTML]{FFFFFF}1}          & 0          & \multicolumn{1}{l|}{\cellcolor[HTML]{ECF4FF}{\color[HTML]{000000} 0}}          & \multicolumn{1}{l|}{\cellcolor[HTML]{ECF4FF}{\color[HTML]{000000} 0}}          & \multicolumn{1}{l|}{\cellcolor[HTML]{ECF4FF}{\color[HTML]{000000} 0}}          & \multicolumn{1}{l|}{\cellcolor[HTML]{ECF4FF}{\color[HTML]{000000} 1}}          & \multicolumn{1}{l|}{\cellcolor[HTML]{ECF4FF}{\color[HTML]{000000} 1}}          & \multicolumn{1}{l|}{\cellcolor[HTML]{ECF4FF}{\color[HTML]{000000} 1}}          & \cellcolor[HTML]{ECF4FF}{\color[HTML]{000000} 1} \\ \hline
{\color[HTML]{9B9B9B} 15}                    & \multicolumn{1}{c|}{\cellcolor[HTML]{FFFFFF}1}          & \multicolumn{1}{c|}{\cellcolor[HTML]{FFFFFF}1}          & \multicolumn{1}{c|}{\cellcolor[HTML]{FFFFFF}1}          & 1          & \multicolumn{1}{l|}{\cellcolor[HTML]{ECF4FF}{\color[HTML]{000000} 0}}          & \multicolumn{1}{l|}{\cellcolor[HTML]{ECF4FF}{\color[HTML]{000000} 0}}          & \multicolumn{1}{l|}{\cellcolor[HTML]{ECF4FF}{\color[HTML]{000000} 0}}          & \multicolumn{1}{l|}{\cellcolor[HTML]{ECF4FF}{\color[HTML]{000000} 0}}          & \multicolumn{1}{l|}{\cellcolor[HTML]{ECF4FF}{\color[HTML]{000000} 0}}          & \multicolumn{1}{l|}{\cellcolor[HTML]{ECF4FF}{\color[HTML]{000000} 0}}          & \cellcolor[HTML]{ECF4FF}{\color[HTML]{000000} 0} \\ \hline
\end{tabular}
\end{table}

\begin{figure}
    \centering
    \includegraphics[width=1.0\linewidth]{img/sim/BCD_7seg_logisim.png}
    \caption{Simulación en logisim del codificador BCD a display de 7 segmentos (ver \href{https://youtu.be/_nBrKwiCBEE}{video de simulación}).}
    \label{fig:BCD_7seg_logisim}
\end{figure}

\begin{figure}
    \centering
    \includegraphics[width=1.0\linewidth]{img/sim/BCD_7seg_proteus.png}
    \caption{Simulación en Proteus del codificador BCD a display de 7 segmentos (ver \href{https://youtu.be/dHLRmH1Clhk?si=OIODbV-4adk1f7gQ}{video de simulación}).}
    \label{fig:BCD_7seg_proteus}
\end{figure}

\begin{figure}
    \centering
    \includegraphics[width=1.0\linewidth]{img/sim/BCD_7seg_ise.png}
    \caption{Sintetización de hardware del codificador BCD a display de 7 segmentos.}
    \label{fig:BCD_7seg_ise}
\end{figure}


\begin{figure}
    \centering
    \includegraphics[width=1.0\linewidth]{img/lab/display_montaje.jpeg}
    \caption{Montaje codificador BCD a 7 segmentos con display de cátodo común (ver \href{https://youtu.be/YC1jDNmOjXA}{video}).}
    \label{fig:BCD_7seg_montaje}
\end{figure}


\subsection{Análisis de resultados}

\begin{itemize}
	\item En cuanto a las compuertas utilizadas, se deben tener las mismas consideraciones mencionadas en la práctica 1 (pull-up o pull-down y tensiones de alimentación de las compuertas).

	\item Para el caso de mintérminos se observó que la simplificación de la salida se reducía a conectarla directamente a la entrada D. Por tanto, las entradas A, B y C correspondían a entradas con un valor de "no importa"~(X).  
	\item Para el caso de maxitérminos (POS) se simplifica la expresión booleana encontrando y agrupando las celdas que contienen $0$ en el mapa de Karnaugh. Y se comprobó que la tabla de verdad se cumple y no afecta que la reducción sea por SOP o POS ya que los circuitos combinaciones son equivalentes por cualquiera de los dos métodos. Las simplificaciones se muestran en la figura~\ref{fig:analisis_0}. La demostración es la siguiente:
	
\begin{equation*}
\begin{matrix} \left ( \bar{C} + \bar{D} \right) \left ( \bar{B} + \bar{D} \right) & = & \bar{C} \bar{B} + \bar{C} \bar{D} + \bar{D} \bar{B} + \bar{D} \bar{D} \\
 & = & \bar{C} \bar{B} + \bar{C} \bar{D} + \bar{D} \bar{B} + \bar{D}\\
 & = & \bar{B} \bar{C} + \bar{D} \left( \bar{C} + \bar{B} + 1 \right) \\
 & = & \bar{B} \bar{C} + \bar{D} \left( 1 \right)\\
 & = & \bar{B} \bar{C} + \bar{D}
\end{matrix}
\end{equation*}	
	
\begin{figure}
    \centering
    \includegraphics[width=1.0\linewidth]{img/sim/analisis_1.png}
    \caption{Simplificación por POS y SOP de la función booleana de la ecuación~\ref{eq:booleana_max}.}
    \label{fig:analisis_0}
\end{figure}

	\item Para esta práctica, al igual que en la primera práctica, se utilizan circuitos combinacionales. En la primera parte son combinaciones de compuertas AND, OR y NOT. Para la segunda parte, debido a su funcionalidad, el integrado 7448 se define como un decodificador de BCD a 7 segmentos de cátodo común. Para éste último se comprobó el funcionamiento esperado con 4 entradas~\footnote{Estas entradas se conectaron resistencia de pull-down para evitar ambigüedades en alta impedancia.} (A, B, C, D) que representaban un número binario, siendo A el bit menos significativo y D el bit más significativo.

	\item Es importante utilizar resistencias limitadoras para alimentar los segmentos del display con las señales $a$, $b$, $c$, $d$, $e$, $f$ y $g$. Cada segmento es un diodo emisor de luz y limitar la corriente que circula sobre ellos es una manera de proteger el elemento y alargar su vida útil.
	
	\item Aunque se utilizaron los diagramas de tiempo para confirmar el comportamiento esperado de las salidas según los valores de las entradas no son de crucial importancia para analizar circuitos combinacionales. Éstos diagramas son muy útiles en el análisis de circuitos secuenciales síncronos y asíncronos.
	
\end{itemize}